\subsection{Sequence Behavior Metrics}

Although a \texttt{range} object does not store all its values, it still behaves like a sequence in several important ways. It has a length, supports indexing, supports slicing, can be traversed in a \texttt{for} loop, and can be tested for membership with \texttt{in}.

The important difference is that these behaviors are computed from the arithmetic descriptor. The \texttt{range} object does not need a stored array of produced integers in order to act like an integer sequence.

\begin{quote}
A \texttt{range} is sequence-like in behavior, but descriptor-like in storage.
\end{quote}

This makes \texttt{range} a good bridge object. It looks simple at the surface, but it already introduces lazy computation, virtual indexing, efficient membership testing, immutability, and reusable traversal.

\subsubsection{Virtual Indexing Lookups: How the \texttt{range} Object Computes Values Arithmetically on Demand}

Indexing a list retrieves a reference from a stored element slot. Indexing a \texttt{range} computes the requested value from the stored \texttt{start}, \texttt{stop}, and \texttt{step} information.

\begin{codebox}[]{python}{Virtual Indexing Lookups: How the range Object Computes Values Arithmetically...}{seq10}
nums = range(10, 20, 2)

print(f"nums: {nums}")
print(f"type(nums): {type(nums)}")
print(f"dir(nums) sample: {dir(nums)[:12]}")

print(f"nums.start: {nums.start}")
print(f"nums.stop: {nums.stop}")
print(f"nums.step: {nums.step}")

print(f"nums[0]: {nums[0]}")
print(f"nums[1]: {nums[1]}")
print(f"nums[2]: {nums[2]}")
print(f"nums[3]: {nums[3]}")
\end{codebox}


\begin{iobox}{Expected Output}
\texttt{nums: range(10, 20, 2)} \\
\texttt{type(nums): <class 'range'>} \\
\texttt{dir(nums) sample: ['\_\_bool\_\_', '\_\_class\_\_', '\_\_contains\_\_', '\_\_delattr\_\_',} \\
\texttt{'\_\_dir\_\_', '\_\_doc\_\_', '\_\_eq\_\_', '\_\_format\_\_', '\_\_ge\_\_',} \\
\texttt{'\_\_getattribute\_\_', '\_\_getitem\_\_', '\_\_gt\_\_']} \\
\texttt{nums.start: 10} \\
\texttt{nums.stop: 20} \\
\texttt{nums.step: 2} \\
\texttt{nums[0]: 10} \\
\texttt{nums[1]: 12} \\
\texttt{nums[2]: 14} \\
\texttt{nums[3]: 16} \\
\end{iobox}


The lookup rule is:

\begin{verbatim}
value = start + index * step
\end{verbatim}

So for \texttt{range(10, 20, 2)}:

\begin{verbatim}
index 0 -> 10 + 0 * 2 -> 10
index 1 -> 10 + 1 * 2 -> 12
index 2 -> 10 + 2 * 2 -> 14
index 3 -> 10 + 3 * 2 -> 16
\end{verbatim}

The \texttt{\_\_getitem\_\_()} method is the special method behind indexing:

\begin{verbatim}
nums = range(10, 20, 2)

print(f"nums[4]: {nums[4]}")
print(f"nums.__getitem__(4): {nums.__getitem__(4)}")
\end{verbatim}

Possible output:

\begin{verbatim}
nums[4]: 18
nums.__getitem__(4): 18
\end{verbatim}

Negative indexing also works, because the range has a computable length:

\begin{verbatim}
nums = range(10, 20, 2)

print(f"len(nums): {len(nums)}")
print(f"nums[-1]: {nums[-1]}")
print(f"nums[-2]: {nums[-2]}")
\end{verbatim}

Possible output:

\begin{verbatim}
len(nums): 5
nums[-1]: 18
nums[-2]: 16
\end{verbatim}

The negative index is translated into a positive position relative to the sequence length:

\begin{verbatim}
nums[-1] -> nums[len(nums) - 1] -> nums[4]
nums[-2] -> nums[len(nums) - 2] -> nums[3]
\end{verbatim}

Out-of-range indexing still fails:

\begin{codebox}[]{python}{Virtual Indexing Lookups: How the range Object Computes Values Arithmetically...}{seq09}
nums = range(10, 20, 2)

try:
    print(nums[99])
except IndexError as err:
    print(f"type(err): {type(err)}")
    print(f"message: {err}")
\end{codebox}


\begin{iobox}{Expected Output}
\texttt{type(err): <class 'IndexError'>} \\
\texttt{message: range object index out of range} \\
\end{iobox}


Slicing a \texttt{range} does not materialize a list. It produces another \texttt{range} object.

\begin{codebox}[]{python}{Virtual Indexing Lookups: How the range Object Computes Values Arithmetically...}{seq08}
nums = range(0, 20, 2)
part = nums[2:7]

print(f"nums: {nums}")
print(f"part: {part}")
print(f"type(part): {type(part)}")
print(f"list(part): {list(part)}")
\end{codebox}


\begin{iobox}{Expected Output}
\texttt{nums: range(0, 20, 2)} \\
\texttt{part: range(4, 14, 2)} \\
\texttt{type(part): <class 'range'>} \\
\texttt{list(part): [4, 6, 8, 10, 12]} \\
\end{iobox}


This confirms the same invariant again: the result of slicing is still descriptive until it is explicitly materialized with \texttt{list()}.

\subsubsection{Arithmetic Membership Testing: Why Integer Containment in \texttt{range} Can Be Checked Without Linear Scanning}

The expression:

\begin{verbatim}
x in range_obj
\end{verbatim}

does not necessarily mean that Python must walk through every value one by one. For integer membership in a \texttt{range}, Python can use arithmetic.

For a value to belong to a positive-step range, three facts must hold:

\begin{itemize}
    \item the value must be greater than or equal to \texttt{start};
    \item the value must be before \texttt{stop};
    \item the distance from \texttt{start} must be exactly divisible by \texttt{step}.
\end{itemize}

Example:

\begin{verbatim}
nums = range(10, 20, 2)

print(f"nums: {nums}")
print(f"12 in nums: {12 in nums}")
print(f"13 in nums: {13 in nums}")
print(f"20 in nums: {20 in nums}")
\end{verbatim}

Possible output:

\begin{verbatim}
nums: range(10, 20, 2)
12 in nums: True
13 in nums: False
20 in nums: False
\end{verbatim}

The value \texttt{12} is present because:

\begin{verbatim}
12 is inside the interval 10 <= value < 20
12 - 10 = 2
2 is divisible by step 2
\end{verbatim}

The value \texttt{13} is not present because:

\begin{verbatim}
13 is inside the interval 10 <= value < 20
13 - 10 = 3
3 is not divisible by step 2
\end{verbatim}

The value \texttt{20} is not present because \texttt{range()} stops before the stop boundary.

The \texttt{\_\_contains\_\_()} method is the special method behind membership testing:

\begin{verbatim}
nums = range(10, 20, 2)

print(f"nums.__contains__(12): {nums.__contains__(12)}")
print(f"nums.__contains__(13): {nums.__contains__(13)}")
\end{verbatim}

Possible output:

\begin{verbatim}
nums.__contains__(12): True
nums.__contains__(13): False
\end{verbatim}

This matters for large ranges:

\begin{verbatim}
nums = range(0, 1_000_000_000, 3)

print(f"nums: {nums}")
print(f"999_999_999 in nums: {999_999_999 in nums}")
print(f"999_999_998 in nums: {999_999_998 in nums}")
\end{verbatim}

Possible output:

\begin{verbatim}
nums: range(0, 1000000000, 3)
999_999_999 in nums: True
999_999_998 in nums: False
\end{verbatim}

Python does not need to generate one billion possible values to answer these questions. For integer containment, it can decide from the arithmetic descriptor.

Negative steps follow the same principle, with the interval direction reversed:

\begin{codebox}[]{python}{Arithmetic Membership Testing: Why Integer Containment in range Can Be Checke...}{seq07}
nums = range(10, 0, -2)

print(f"nums: {nums}")
print(f"list(nums): {list(nums)}")
print(f"8 in nums: {8 in nums}")
print(f"7 in nums: {7 in nums}")
print(f"0 in nums: {0 in nums}")
\end{codebox}


\begin{iobox}{Expected Output}
\texttt{nums: range(10, 0, -2)} \\
\texttt{list(nums): [10, 8, 6, 4, 2]} \\
\texttt{8 in nums: True} \\
\texttt{7 in nums: False} \\
\texttt{0 in nums: False} \\
\end{iobox}


Again, the stop boundary is excluded. The value \texttt{0} is the boundary, not a produced element.

The important control-flow implication is that \texttt{range} is not only memory-efficient. It also supports efficient decision checks. A condition such as this:

\begin{verbatim}
n = 42

if n in range(0, 100, 2):
    print("n is an accepted even value")
else:
    print("n is outside the accepted sequence")
\end{verbatim}

Possible output:

\begin{verbatim}
n is an accepted even value
\end{verbatim}

is not equivalent to manually scanning a list of values. For integer membership, the range can answer the containment question arithmetically.

\subsubsection{Immutability and Reusability: Using a Single Range Descriptor Across Multiple Traversal Pipelines}

A \texttt{range} object is immutable. Its \texttt{start}, \texttt{stop}, and \texttt{step} attributes cannot be changed after construction, and its virtual elements cannot be assigned.

\begin{codebox}[]{python}{Immutability and Reusability: Using a Single Range Descriptor Across Multiple...}{seq06}
nums = range(0, 5)

print(f"nums: {nums}")
print(f"type(nums): {type(nums)}")
print(f"nums.start: {nums.start}")
print(f"nums.stop: {nums.stop}")
print(f"nums.step: {nums.step}")
\end{codebox}


\begin{iobox}{Expected Output}
\texttt{nums: range(0, 5)} \\
\texttt{type(nums): <class 'range'>} \\
\texttt{nums.start: 0} \\
\texttt{nums.stop: 5} \\
\texttt{nums.step: 1} \\
\end{iobox}


Trying to change an element fails:

\begin{codebox}[]{python}{Immutability and Reusability: Using a Single Range Descriptor Across Multiple...}{seq05}
nums = range(0, 5)

try:
    nums[0] = 99
except TypeError as err:
    print(f"type(err): {type(err)}")
    print(f"message: {err}")
\end{codebox}


\begin{iobox}{Expected Output}
\texttt{type(err): <class 'TypeError'>} \\
\texttt{message: 'range' object does not support item assignment} \\
\end{iobox}


Trying to change an attribute also fails:

\begin{codebox}[]{python}{Immutability and Reusability: Using a Single Range Descriptor Across Multiple...}{seq04}
nums = range(0, 5)

try:
    nums.start = 100
except AttributeError as err:
    print(f"type(err): {type(err)}")
    print(f"message: {err}")
\end{codebox}


\begin{iobox}{Expected Output}
\texttt{type(err): <class 'AttributeError'>} \\
\texttt{message: readonly attribute} \\
\end{iobox}


The range object has no list-like mutation methods such as \texttt{append()}:

\begin{verbatim}
nums = range(0, 5)

print(f"has append: {'append' in dir(nums)}")
print(f"has __getitem__: {'__getitem__' in dir(nums)}")
print(f"has __contains__: {'__contains__' in dir(nums)}")
print(f"has __iter__: {'__iter__' in dir(nums)}")
\end{verbatim}

Possible output:

\begin{verbatim}
has append: False
has __getitem__: True
has __contains__: True
has __iter__: True
\end{verbatim}

So a \texttt{range} is not a mutable container. It is an immutable arithmetic descriptor with sequence behavior.

This immutability supports safe reuse. The same \texttt{range} object can be traversed several times:

\begin{codebox}[]{python}{Immutability and Reusability: Using a Single Range Descriptor Across Multiple...}{seq03}
nums = range(1, 4)

print("first traversal")
for n in nums:
    print(f"n: {n}")

print("second traversal")
for n in nums:
    print(f"n: {n}")
\end{codebox}


\begin{iobox}{Expected Output}
\texttt{first traversal} \\
\texttt{n: 1} \\
\texttt{n: 2} \\
\texttt{n: 3} \\
\texttt{second traversal} \\
\texttt{n: 1} \\
\texttt{n: 2} \\
\texttt{n: 3} \\
\end{iobox}


This works because the \texttt{range} object is iterable, not an already-consumed iterator. Each new \texttt{for} loop asks the range object for a fresh iterator.

This can be shown explicitly:

\begin{codebox}[]{python}{Immutability and Reusability: Using a Single Range Descriptor Across Multiple...}{seq02}
nums = range(1, 4)

it1 = iter(nums)
it2 = iter(nums)

print(f"type(it1): {type(it1)}")
print(f"type(it2): {type(it2)}")
print(f"it1 is it2: {it1 is it2}")

print(f"next(it1): {next(it1)}")
print(f"next(it1): {next(it1)}")

print(f"next(it2): {next(it2)}")
\end{codebox}


\begin{iobox}{Expected Output}
\texttt{type(it1): <class 'range\_iterator'>} \\
\texttt{type(it2): <class 'range\_iterator'>} \\
\texttt{it1 is it2: False} \\
\texttt{next(it1): 1} \\
\texttt{next(it1): 2} \\
\texttt{next(it2): 1} \\
\end{iobox}


The two iterators are separate traversal state machines over the same immutable descriptor. Advancing \texttt{it1} does not advance \texttt{it2}.

This distinction will matter later when iterator exhaustion is discussed. For now, the key difference is:

\begin{itemize}
    \item \texttt{range} is reusable because it can produce new iterators.
    \item a specific \texttt{range\_iterator} is consumed as it advances.
\end{itemize}

A range can therefore be shared across several traversal pipelines:

\begin{codebox}[]{python}{Immutability and Reusability: Using a Single Range Descriptor Across Multiple...}{seq01}
steps = range(1, 6)

print("plain traversal")
for n in steps:
    print(f"n: {n}")

print("squared values")
for n in steps:
    print(f"{n} squared is {n * n}")

print("formatted table")
for n in steps:
    print(f"{n:>2} | {n * 42:>3}")
\end{codebox}


\begin{iobox}{Expected Output}
\texttt{plain traversal} \\
\texttt{n: 1} \\
\texttt{n: 2} \\
\texttt{n: 3} \\
\texttt{n: 4} \\
\texttt{n: 5} \\
\texttt{squared values} \\
\texttt{1 squared is 1} \\
\texttt{2 squared is 4} \\
\texttt{3 squared is 9} \\
\texttt{4 squared is 16} \\
\texttt{5 squared is 25} \\
\texttt{formatted table} \\
\texttt{ 1 |  42} \\
\texttt{ 2 |  84} \\
\texttt{ 3 | 126} \\
\texttt{ 4 | 168} \\
\texttt{ 5 | 210} \\
\end{iobox}


No traversal changes the \texttt{steps} object. The descriptor remains the same.

The practical rule is:

\begin{quote}
A \texttt{range} object is a stable, immutable description of an arithmetic sequence. Traversing it creates iterator state, but the range object itself is not consumed.
\end{quote}

This is why \texttt{range} is so useful in control-flow code. It supplies repeatable arithmetic traversal without materializing a list and without changing itself during traversal.